%% appendix_notation.tex — Acronyms, Notation, and Data Source Details
%% Extracted from main.tex
\chapter{Acronyms, Notation, and Data Source Details}
\label{app:notation}


\section{Acronyms and Nomenclature}
\label{sec:acronyms}

The following acronyms are used throughout this chapter and the remainder of
the thesis.  They are defined here for convenience and also at the point of
first use.

\begin{description}[labelwidth=3.0cm,leftmargin=3.5cm]
  \item[SAR]  Synthetic Aperture Radar.  An active microwave imaging technique
    in which the motion of the spacecraft is used to synthesise a very long
    antenna aperture, yielding fine-resolution imagery of surface texture and
    composition.
  \item[VIMS] Visible and Infrared Mapping Spectrometer.  The \cassini{}
    instrument that imaged \titan{}'s surface through six near-infrared
    atmospheric transparency windows, providing surface composition data.
  \item[CIRS] Composite Infrared Spectrometer.  The \cassini{} thermal
    emission instrument used here to constrain the surface temperature field.
  \item[GTDR] Global Titan Digital elevation model, Radar --- the raw
    Cassini RADAR altimetry dataset archived in the Planetary Data System.
  \item[GTDE] Global Titan Digital Elevation model --- the globally
    interpolated topographic product of \citet{Corlies2017Titan} derived from
    the GTDR and supplementary data sources.
  \item[DEM]  Digital Elevation Model.  A gridded representation of surface
    topography.
  \item[PDS]  Planetary Data System.  The NASA archive for planetary mission
    data.
  \item[HDI]  Highest Density Interval.  A Bayesian credible interval
    containing the specified probability mass within the narrowest possible
    range.
\end{description}

\section{Notation}
\label{sec:notation}

Two quantities appear throughout this chapter and are defined here to avoid
repetition.  Each habitability-proxy feature $f_i$ ($i = 1,\ldots,8$) is
assigned:

\begin{itemize}
  \item A \textbf{weight} $w_i \in [0,1]$, where $\sum_i w_i = 1$.  The weight
    encodes the relative astrobiological importance of feature $i$ as a
    habitability indicator.  A feature with a higher weight exerts a larger
    influence on the posterior probability of habitability, all else being
    equal.  Weights are assigned from physical reasoning and validated against
    data-derived importances (Section~\ref{sec:importances}).

  \item A \textbf{prior mean} $\mu_i \in [0,1]$.  This is the expected
    spatially-averaged value of feature $i$ before any observational data
    are considered --- a statement of our prior knowledge about how common or
    widespread the habitability-relevant condition is on \titan{}.  For
    example, liquid hydrocarbon covers $\sim$2\% of \titan{}'s surface
    \citep{Hayes2008}, so its prior mean is $\mu_1 = 0.02$.  The weighted
    sum of all prior means defines the global prior mean habitability:
    $\mu_0 = \sum_i w_i \mu_i = 0.331$.  Note that the value $\mu_0 = 0.331$ does not appear in the prior literature as a global habitability estimate for \titan{}; it emerges from the component prior means in Table~\ref{tab:priors} and should be interpreted as encoding the observation that \titan{}'s surface is globally organic-rich whilst simultaneously recognising that liquid solvent is present at only $\sim$2\% of the surface. The resulting moderate prior is intentionally revisionable: the weak concentration parameter $\kappa = 5$ ensures that the data, not the prior, determine the posterior at any pixel where observational coverage is good. Full values are tabulated in
    Appendix~\ref{app:priors}.
\end{itemize}

This notation first appears in a practical context in Section~\ref{sec:features},
where the eight features are described individually.  Their role in the
Bayesian update is explained in Section~\ref{sec:bayes}.


\section{Observational Data Source Details}
\label{app:data_sources}
\label{sec:data_sources}

\subsection{Cassini RADAR Synthetic Aperture Radar Mosaic}
\label{sec:sar_data}

The Cassini RADAR instrument operated in SAR mode during \titan{} flybys,
imaging the surface at $\sim$350\,m per pixel over the 127 targeted flybys
from 2004 to 2017 \citep{Janssen2016}.  SAR measures the normalised radar
cross-section $\sigma^0$ (typically expressed in dB), which is sensitive to
surface roughness, dielectric properties, and the presence of smooth or
conducting materials.  Low backscatter ($\sigma^0 < -20$\,dB) is
diagnostic of smooth surfaces (consistent with liquid hydrocarbon lakes) or
dielectrically absorbing material (consistent with tholin-rich dune sand).
High backscatter correlates with rough, water-ice-rich material such as
mountain terrain, crater ejecta, and some cryovolcanic surfaces.  The SAR
mosaic covers $\sim$67\% of \titan{}'s surface and constitutes the primary
data source for the liquid hydrocarbon, chemical energy, and subsurface ocean
features.

\subsection{VIMS Near-Infrared Spectral Mosaic}
\label{sec:vims_data}

The Cassini Visible and Infrared Mapping Spectrometer (VIMS) observed \titan{}'s
surface through atmospheric transparency windows at 0.94, 1.08, 1.27, 1.59,
2.03, and 2.79\,\textmu{}m.  The band ratio 1.59/1.27\,\textmu{}m is an
established proxy for tholin (complex organic) abundance at the surface: the
1.59\,\textmu{}m window probes deeper into the tholin-bearing surface whilst
the 1.27\,\textmu{}m window provides atmospheric normalisation \citep{Barnes2007}.
The Seignovert et al.\ (2019) combined VIMS+ISS mosaic \citep{Seignovert2019CaltechDATA}
provides this band ratio map at 2--4\,km per pixel coverage, but only
for the 0--180\textdegree{}W hemisphere ($\sim$50\% of the globe).  The
significance of this coverage gap for the organic abundance feature is
discussed in Section~\ref{sec:feat_organic}.

\subsection{CIRS Thermal Emission Observations}
\label{sec:cirs_data}

The Composite Infrared Spectrometer (CIRS) measured \titan{}'s surface
brightness temperature at far-infrared wavelengths throughout the Cassini
mission.  \citet{Jennings2019} derived an analytical model of the
latitude-dependent surface temperature from 13 years of CIRS observations,
producing the formula:
\begin{equation}
  T(\phi, Y) = (93.53 - 0.095\,Y)\cos\!\Bigl[
    (\phi + 0.85 - 3.2\,Y)(0.0029 - 0.00006\,Y)
  \Bigr]
  \label{eq:jennings}
\end{equation}
where $\phi$ is latitude in degrees and $Y = t_\text{CE} - 2009.61$ is years
from the northern spring equinox (22 May 2009).  This formula reproduces the
observed zonal surface temperatures to within $\pm0.4$\,K (rms) across the
full Cassini epoch and provides 100\% global coverage, grounded in CIRS
focal-plane-1 limb and nadir observations across the full 2004--2017 mission.
The model captures only latitudinal and seasonal variation; no longitudinal
inhomogeneity is represented (consistent with the absence of significant
longitudinal temperature gradients in the Cassini dataset).  For temporal
epochs other than the present, the mean annual temperature field is held
stationary: this approximation is justified because the dominant habitability
driver at non-present epochs is the liquid-hydrocarbon scale function
$s_1(t)$ rather than the temperature gradient, and the methane-cycle
feature contributes only $w_4 = 0.15$ to the posterior.  It is used to
synthesise the global temperature field and its meridional gradient, which
enters the methane cycle feature (Section~\ref{sec:feat_methane}).

\subsection{GTDE/GTDR Topographic Models}
\label{sec:topography_data}

The Global Titan Digital Elevation model (GTDE), produced by
\citet{Corlies2017Titan}, integrates Cassini RADAR altimetry, SARtopo
stereophotogrammetry, and radial basis function interpolation to produce
a globally-interpolated topographic model at $\sim$22\,km per pixel (2 pixels
per degree).  Coverage is approximately 90\% of the globe, with the south
polar region and some high-latitude northern areas relying most heavily on
interpolation.  The topographic model contributes to the chemical energy
feature (identifying topographic lows as preferential organic accumulation
sites), the topographic complexity feature (through local roughness estimation),
and the surface--atmosphere interaction feature (through slope calculation).

\subsection{Lopes (2019) Global Geomorphological Map}
\label{sec:geo_data}

The global geomorphological map of \citet{Lopes2019} classifies the Cassini
SAR-imaged surface into seven terrain units (plains, dunes, mountains, basins,
labyrinth terrain, craters, and lakes) using SAR texture, backscatter,
morphology, and VIMS spectral correlation.  This map provides the terrain
classification that enters the organic abundance feature (as a globally
seamless gap-fill) and the geomorphological diversity feature.  The shapefiles
are provided in east-positive geographic coordinates and require conversion to
the west-positive convention of the analysis grid.

\subsection{Miller (2021) Global Fluvial Channel Map}
\label{sec:channel_data}

\citet{Miller2021channels} identified and mapped \titan{}'s global fluvial
channel network from Cassini SAR data, producing a vector dataset of channel
locations and widths.  This network constitutes the primary mechanism for
organic material transport from equatorial source regions to the polar basins,
and channel density is used as one component of the surface--atmosphere
interaction feature.

\subsection{Tidal Love Number and Gravity Data}
\label{sec:gravity_data}

The degree-2 tidal Love number $k_2 = 0.589 \pm 0.150$ from \citet{Iess2012}
provides a global constraint on interior structure, confirming the presence
of a subsurface liquid layer.  This scalar quantity enters as a global floor
value in the subsurface ocean proximity feature.

\subsection{Hedgepeth (2020) Global Crater Catalogue}
\label{sec:crater_data}

The global catalogue of confirmed and probable impact structures on \titan{}
by \citet{Hedgepeth2020} records crater diameters and positions for all
resolved impact features in the Cassini SAR dataset.  This catalogue is used
exclusively in the \textit{Past} epoch to generate the impact-melt proximity feature.
